% Part: first-order-logic
% Chapter: completeness
% Section: maximally-consistent-sets

% Definition of maximally consistent sets. Properties of mcs required
% for completeness proved are in maximally-consistent-sets.tex in the
% chapter on the proof system used.

\documentclass[../../../include/open-logic-section]{subfiles}

\begin{document}

\olfileid{fol}{com}{mcs}
\olsection{\usetoken{P}{sentence} యొక్క గరిష్ఠ అవైరుధ్య సమితులు}

\begin{defn}[గరిష్ఠ అవైరుధ్య సమితి]
\ollabel{mcs}
\tetoken{వాక్యాల}{!!{sentence}s} సమితి~$\Gamma$ \emph{గరిష్ఠ అవైరుధ్యమైనది} కావడానికి, కావలసినవీ చాలునవీ:
\begin{enumerate}
\item $\Gamma$ అవైరుధ్యమైనది, మరియు
\item $\Gamma \subsetneq \Gamma'$ అయితే $\Gamma'$ వైరుధ్యపూరితం.
\end{enumerate}
\end{defn}

\begin{explain}
పైనిదానికి తుల్యమైన ప్రత్యామ్నాయ నిర్వచనం: వాక్యాల సమితి $\Gamma$ \emph{గరిష్ఠ అవైరుధ్యమైనది} కావడానికి, కావలసినవీ చాలునవీ:
\begin{enumerate}
\item $\Gamma$ అవైరుధ్యమైనది, మరియు
\item $\Gamma \cup \{ !A \}$ అవైరుధ్యమైతే $!A \in \Gamma$.
\end{enumerate}
మరో మాటలో, $\Gamma$లో అప్పటికే లేని \tetoken{వాక్యాలను}{!!{sentence}s} గరిష్ఠ అవైరుధ్య సమితి~$\Gamma$కి జోడించి, ఫలితంగా వచ్చే పెద్ద సమితిని అవైరుధ్యంగానే ఉంచలేం.
\end{explain}

\begin{explain}
సంపూర్ణత నిరూపణలో గరిష్ఠ అవైరుధ్య సమితులు ముఖ్యమైనవి. ఎందుకంటే ప్రతి అవైరుధ్య \tetoken{వాక్యాల}{!!{sentence}s} సమితి~$\Gamma$, ఏదో గరిష్ఠ అవైరుధ్య సమితి~$\Gamma^*$లో ఇమిడి ఉంటుందని హామీ ఇవ్వగలం. అలాగే గరిష్ఠ అవైరుధ్య సమితిలో ప్రతి \tetoken{వాక్యానికి}{!!{sentence}}~$!A$, $!A$ లేదా దాని నిరాకరణ $\lnot !A$ ఉంటుంది. ముఖ్యంగా అణు \tetoken{వాక్యాలకూ}{!!{sentence}s} ఇది వర్తిస్తుంది. కనుక \tetoken{స్థిరాంకాలతో}{!!{constant}s} తగినట్లు విస్తరించిన భాషలోని గరిష్ఠ అవైరుధ్య సమితి నుంచి \tetoken{ఒక నిర్మాణాన్ని}{!!a{structure}} నిర్మించవచ్చు; అందులో \tetoken{విధేయాల}{!!{predicate}s} అన్వయాన్ని $\Gamma^*$లో ఏ అణు \tetoken{వాక్యాలు}{!!{sentence}s} ఉన్నాయనే దాని ప్రకారం నిర్వచిస్తాం. తరువాత ఈ \tetoken{నిర్మాణం}{!!{structure}}, $\Gamma^*$లోని అన్ని \tetoken{వాక్యాలను}{!!{sentence}s} (అందువల్ల $\Gamma$లోని వాటిని కూడా) సత్యం చేస్తుందని చూపవచ్చు. ఈ చివరి విషయ నిరూపణకు $\lnot !A \in
\Gamma^*$ కావడానికి, కావలసినదీ చాలునదీ $!A \notin \Gamma^*$ కావడం; అలాగే $(!A \lor !B) \in \Gamma^*$ కావడానికి, కావలసినదీ చాలునదీ $!A
\in \Gamma^*$ లేదా $!B \in \Gamma^*$ కావడం మొదలైనవి అవసరం.
\end{explain}

\begin{prop}
\ollabel{prop:mcs}
$\Gamma$ గరిష్ఠ అవైరుధ్యమైనదని అనుకుందాం. అప్పుడు:
\begin{enumerate}
\item \ollabel{prop:mcs-prov-in} $\Gamma \Proves !A$ అయితే $!A \in
  \Gamma$.

\item \ollabel{prop:mcs-either-or} ఏ $!A$కైనా $!A \in
  \Gamma$ లేదా $\lnot !A \in \Gamma$.

\tagitem{prvAnd}{\ollabel{prop:mcs-and} $(!A \land !B) \in \Gamma$
  కావడానికి, కావలసినదీ చాలునదీ $!A \in \Gamma$, $!B \in \Gamma$ రెండూ.}{}

\tagitem{prvOr}{\ollabel{prop:mcs-or} $(!A \lor !B) \in \Gamma$ కావడానికి, కావలసినదీ చాలునదీ $!A \in \Gamma$ లేదా $!B \in \Gamma$.}{}

\tagitem{prvIf}{\ollabel{prop:mcs-if} $(!A \lif !B) \in \Gamma$ కావడానికి, కావలసినదీ చాలునదీ $!A \notin \Gamma$ లేదా $!B \in \Gamma$.}{}
\end{enumerate}
\end{prop}

\begin{proof}
క్రింది వాటన్నింటిలో $\Gamma$ గరిష్ఠ అవైరుధ్యమైనదని అనుకుందాం.
\begin{enumerate}
\item $\Gamma \Proves !A$ అయితే $!A \in \Gamma$.

$\Gamma \Proves !A$ అనుకుందాం. దీనికి విరుద్ధంగా $!A
\notin \Gamma$ అనుకుందాం. $\Gamma$ గరిష్ఠ అవైరుధ్యమైనది కనుక $\Gamma
\cup \{!A\}$ వైరుధ్యపూరితం; కాబట్టి $\Gamma \cup \{!A\} \Proves
\lfalse$. \tagrefs{prfSC/{fol:seq:prv:prop:provability-contr},prfND/{fol:ntd:prv:prop:provability-contr}} ప్రకారం $\Gamma$ వైరుధ్యపూరితం. కానీ ఇది $\Gamma$ అవైరుధ్యమనే పరికల్పనకు విరుద్ధం. అందువల్ల $!A \notin
\Gamma$ అనడం కుదరదు; కాబట్టి $!A \in \Gamma$.

\item ఏ $!A$కైనా $!A \in \Gamma$ లేదా $\lnot !A \in \Gamma$.

విరుద్ధంగా, ఏదో~$!A$కు $!A \notin \Gamma$, $\lnot !A \notin \Gamma$ రెండూ అనుకుందాం. $\Gamma$ గరిష్ఠ అవైరుధ్యమైనది కనుక $\Gamma \cup \{!A\}$, $\Gamma \cup \{\lnot !A\}$ రెండూ వైరుధ్యపూరితాలు. అంటే $\Gamma \cup \{!A\} \Proves \lfalse$, $\Gamma \cup
\{\lnot !A\} \Proves \lfalse$. \tagrefs{prfSC/{fol:seq:prv:prop:provability-exhaustive},prfND/{fol:ntd:prv:prop:provability-exhaustive}} ప్రకారం $\Gamma$ వైరుధ్యపూరితం అవుతుంది---ఇది వైరుధ్యం. కాబట్టి అటువంటి \tetoken{వాక్యం}{!!a{sentence}}~$!A$ ఉండదు; ప్రతి $!A$కు $!A \in \Gamma$ లేదా $\lnot !A
\in \Gamma$.

\tagitem{defAnd}{}{%
\iftag{probAnd}{సాధన.}{%
$(!A \land !B) \in \Gamma$ కావడానికి, కావలసినదీ చాలునదీ $!A \in \Gamma$, $!B \in \Gamma$ రెండూ:

ముందు దిశకు $(!A \land !B) \in \Gamma$ అనుకుందాం. అప్పుడు $\Gamma \Proves !A \land !B$. \tagrefs{prfSC/{fol:seq:prv:prop:provability-land-left},prfND/{fol:ntd:prv:prop:provability-land-left}} ప్రకారం $\Gamma \Proves !A$, $\Gamma \Proves !B$. \olref{prop:mcs-prov-in} ప్రకారం కావలసినట్లు $!A \in \Gamma$, $!B \in \Gamma$.

వెనుక దిశకు $!A \in \Gamma$, $!B \in
\Gamma$ అనుకుందాం. అప్పుడు $\Gamma \Proves !A$, $\Gamma \Proves !B$. \tagrefs{prfSC/{fol:seq:prv:prop:provability-land-right},prfND/{fol:ntd:prv:prop:provability-land-right}} ప్రకారం $\Gamma \Proves !A \land !B$. \olref{prop:mcs-prov-in} ప్రకారం $(!A \land
!B) \in \Gamma$.}}

\tagitem{defOr}{}{%
\iftag{probOr}{సాధన.}{%
$(!A \lor !B) \in \Gamma$ కావడానికి, కావలసినదీ చాలునదీ $!A \in \Gamma$ లేదా $!B \in \Gamma$.

ముందు దిశకు వ్యతిరేకాన్వయంగా $!A
\notin \Gamma$, $!B \notin \Gamma$ అనుకుందాం. $(!A \lor
!B) \notin \Gamma$ అని చూపాలి. $\Gamma$ గరిష్ఠ అవైరుధ్యమైనది కనుక $\Gamma
\cup \{!A\} \Proves \lfalse$, $\Gamma \cup \{!B\} \Proves \lfalse$. \tagrefs{prfSC/{fol:seq:prv:prop:provability-lor-left},prfND/{fol:ntd:prv:prop:provability-lor-left}} ప్రకారం $\Gamma \cup \{(!A \lor !B)\}$ వైరుధ్యపూరితం. అందుచేత కావలసినట్లు $(!A \lor !B)
\notin \Gamma$.

వెనుక దిశకు $!A \in \Gamma$ లేదా $!B \in
\Gamma$ అనుకుందాం. అప్పుడు $\Gamma \Proves !A$ లేదా $\Gamma \Proves !B$. \tagrefs{prfSC/{fol:seq:prv:prop:provability-lor-right},prfND/{fol:ntd:prv:prop:provability-lor-right}} ప్రకారం $\Gamma \Proves !A \lor !B$. \olref{prop:mcs-prov-in} ప్రకారం కావలసినట్లు $(!A \lor
!B) \in \Gamma$.}}

\tagitem{defIf}{}{%
\iftag{probIf}{సాధన.}{%
$(!A \lif !B) \in \Gamma$ కావడానికి, కావలసినదీ చాలునదీ $!A \notin \Gamma$ లేదా $!B \in \Gamma$:

ముందు దిశకు $(!A \lif !B) \in \Gamma$ అనుకుందాం; దీనికి విరుద్ధంగా $!A \in \Gamma$, $!B \notin \Gamma$ అనుకుందాం. ఈ పరికల్పనల ప్రకారం $\Gamma \Proves !A \lif !B$, $\Gamma \Proves !A$. \tagrefs{prfSC/{fol:seq:prv:prop:provability-mp},prfND/{fol:ntd:prv:prop:provability-mp}} ప్రకారం $\Gamma \Proves !B$. కానీ \olref{prop:mcs-prov-in} వల్ల $!B \in
\Gamma$; ఇది $!B \notin \Gamma$ అనే పరికల్పనకు విరుద్ధం.

వెనుక దిశకు, ముందుగా $!A \notin
\Gamma$ సందర్భాన్ని చూద్దాం. \olref{prop:mcs-either-or} ప్రకారం $\lnot !A \in \Gamma$; అందుచేత $\Gamma \Proves \lnot !A$. \tagrefs{prfSC/{fol:seq:prv:prop:provability-lif},prfND/{fol:ntd:prv:prop:provability-lif}} ప్రకారం $\Gamma \Proves !A \lif !B$. మళ్లీ \olref{prop:mcs-prov-in} వల్ల కావలసినట్లు $(!A \lif !B) \in \Gamma$.

ఇప్పుడు $!B \in \Gamma$ సందర్భాన్ని చూద్దాం. అప్పుడు $\Gamma \Proves !B$; \tagrefs{prfSC/{fol:seq:prv:prop:provability-lif},prfND/{fol:ntd:prv:prop:provability-lif}} ప్రకారం $\Gamma \Proves !A \lif !B$. \olref{prop:mcs-prov-in} వల్ల $(!A \lif
!B) \in \Gamma$.}}
\end{enumerate}
\end{proof}

\begin{prob}
\olref[fol][com][mcs]{prop:mcs} నిరూపణను పూర్తి చేయండి.
\end{prob}

\end{document}
